\chapter{Experimental Setup}

The report separates nominal evaluation, historical PVT qualification, and
optimizer comparison. Numbers from different candidates or runs are not merged.

\section{Recorded v3 training run}

The preserved \code{results/v3\_tt\_1000\_summary.json} records a real-backend
v3 run with seed 101, 100 policy updates, and 692 evaluations. Its best recorded
reward was 10.3096. This artifact establishes that a real-SPICE training run
completed and produced a checkpoint; it does not by itself prove global
convergence or optimizer superiority.

\section{Completed reduced 27-condition qualification}

The primary artifact is
\code{results/web\_ui\_runs/07acbcab1cd346cca24eab31919496c0.json}.
It evaluates selected design \code{pipeline\_ep0} over TT, SS, and FF; three
supply values; and temperatures of 0, 27, and \SI{125}{\degreeCelsius}. All 27
conditions passed at candidate fidelity using a 512-bit request; every condition
records 425 useful simulated bits after alignment and warm-up. The channel is
\code{ieee802\_ibm\_20db\_thru.s4p}, checksum
\code{ce6002f8663623eb...}, with port map 1,3,2,4. This is a completed reduced
grid, not the full five-corner competition matrix.

\section{Fair optimizer comparison}

The no-warm-start comparison gives each method 20 evaluations with seed 123 and
the same success criterion. Random Search samples globally in the continuous
space; CEM uses a centered Gaussian; PPO begins at the parameter-grid midpoint
and uses horizon-one local actions. The budget is matched, but these action-space
differences remain.

\section{Experimental identity}

For any new run, the following fields should accompany the reported result:
repository revision and dirty status; Python, ngspice, and model versions;
design identifier and full parameters; target; channel checksum and port map;
seed and budget; fidelity; PVT set; and result-artifact checksum. Appendix
\ref{app:reproducibility} gives the corresponding commands.

\section{Verification strategy}

The automated suite covers parameter parsing and bounds, netlist rendering,
ngspice process/error handling, channel validation, receiver metrics, PPO
state/action/reward schemas, checkpoint compatibility, event emission, caching,
baseline wiring, PVT selection, specification generation, schematic export, run
graphs, and web-command construction. Optional integration tests exercise
ngspice and SKY130 only when external dependencies are configured.

Passing tests establish software behavior for the tested cases; they do not
establish analog performance over unsampled parameters or conditions. The final
submission should include a dated log from a clean commit and separately report
unit, simulator-integration, and skipped tests.

The most useful end-to-end invariants are that all eight v3 parameters reach
every relevant bench and artifact; cache identity changes with every electrical
input; missing metrics stay missing; incompatible checkpoints fail before
simulation; and nominal success is never promoted to PVT success.
